%==============================================================================
%  25-ai-use-playbook.tex — AI Use Playbook by Function
%==============================================================================
\strategyPart{25 \textbar\ Use Playbook}{AI Use Playbook by Function}

This playbook condenses how AI is applied across business functions, with the
measured benefits and the controls each function requires.

\section{Customer experience}

\textbf{Chatbots} provide 24-hour service with lower cost per interaction;
\textbf{copilots} support agents by searching history, suggesting responses
and detecting sentiment — cutting average handling time and accelerating
training; \textbf{personalization} recommends products, content and
next-best actions; \textbf{churn prediction} flags at-risk customers;
\textbf{voice-of-customer} analysis mines surveys, tickets and transcripts
[WEF; MIT/NBER 2023].

\section{Marketing \& sales}

\textbf{Market analysis} monitors competitors and pricing; \textbf{segmentation}
groups customers by behavior and profitability; \textbf{campaign creation}
generates copy, images and variants (with human review for accuracy, brand
and legal); \textbf{lead scoring} ranks prospects; \textbf{sales forecasting}
flags optimistic or at-risk pipelines. Measured impact: \textbf{55.8\% faster}
coding in engineering and \textbf{+14\%} support productivity translate
directly to marketing/sales automation gains [GitHub 2022; MIT/NBER 2023].

\section{Product management}

AI supports \textbf{customer discovery} (interview transcripts, tickets,
reviews), \textbf{opportunity analysis}, \textbf{requirements development},
\textbf{roadmap prioritization}, \textbf{product analytics} and
\textbf{AI-native product design} — conversational analytics, autonomous
workflow platforms and intelligent search. The product manager remains
responsible for the decision and business context.

\section{Business intelligence \& decision-making}

\textbf{Natural-language data analysis} ("why did revenue drop?"), automated
reporting, anomaly detection, predictive analytics and \textbf{multi-agent BI}
— a set of specialized agents (discovery, query, analysis, visualization,
validation, governance) that make analytics accessible to non-technical users.
Controls: access permissions, query validation, source citation, human
approval for high-impact decisions, audit logs.

\section{Finance, operations, supply chain}

\textbf{Finance:} forecasting, invoice/document extraction, fraud detection,
draft reporting, risk analysis. \textbf{Operations:} process automation,
intelligent process management, predictive maintenance, quality control,
intelligent scheduling. \textbf{Supply chain:} demand forecasting, inventory
optimization, supplier-risk monitoring, logistics optimization and scenario
planning. For regulated financial decisions, AI supports rather than replaces
accountable humans [EU AI Act; NIST AI RMF].

\section{HR, R\&D, IT, legal}

\textbf{HR:} recruitment support, workforce planning, learning
recommendations, employee support assistants and skills-based organization
— with strict controls against discrimination. \textbf{R\&D:} research
discovery, generative design, experiment prioritization, software
development and time-to-market reduction. \textbf{IT/security:} ticket
classification, monitoring, anomaly and threat detection. \textbf{Legal:}
contract analysis, regulatory monitoring and compliance testing.

\section{AI as an autonomous workflow agent}

An agent interprets a goal, retrieves data, plans actions, uses approved
tools, validates results, requests approval where necessary, executes and
logs. Every agent needs defined \textbf{identity, role, permissions, tool
scope, data scope, spending limits, escalation paths and logging} — Deloitte
reports organizations increasingly customize agents, making identity,
authorization, monitoring and human accountability essential
[Deloitte 2026; Gartner 2024].
